\chapter{Microkernel dan Plugin}
\label{bab:plugin}

\section*{Tujuan Pembelajaran}

\begin{enumerate}
  \item Menjelaskan pembagian antara \textit{core} dan \textit{plugin} beserta \textit{contract}
    yang menghubungkan keduanya.
  \item Membedakan \textit{eager activation} dari \textit{lazy activation}, lalu menentukan kapan
    masing-masing bentuk layak dipilih.
  \item Mengukur \textit{core independence}, \textit{error isolation}, dan
    \textit{activation cost}.
\end{enumerate}

\section{Pendahuluan}

Microkernel memisahkan sebuah sistem menjadi \textit{core} yang kecil beserta sejumlah
bagian tambahan yang dapat dipasang dan dicabut. Istilah tersebut lahir di
dunia sistem operasi, yaitu ketika penjadwalan dan pengelolaan memori
dipertahankan di dalam kernel, sedangkan sistem \textit{file} dan pengandar perangkat
dipindahkan ke luar sebagai proses tersendiri.

Bab ini hanya menyinggung sistem operasi sekilas, sebab bentuk yang jauh lebih
sering ditemui mahasiswa adalah bentuk turunannya, yaitu arsitektur
\textit{plugin}. Eclipse, Visual Studio Code, Chrome, dan Firefox memakai bentuk
tersebut. Pada keempatnya, \textit{core} yang dipasang pengguna sebenarnya kecil, dan
hampir seluruh kemampuan yang dipakai sehari-hari datang dari bagian tambahan.

Pertanyaan yang dijawab bab ini karena itu bukan bagaimana menulis kernel,
melainkan bagaimana sebuah \textit{core} dapat menerima kemampuan baru tanpa mengenal
satu pun penyumbangnya. Pembahasan mengikuti kerangka yang disusun oleh Richards
dan Ford~\cite{richards2020fundamentals}, serta catatan Bolour tentang
arsitektur \textit{plugin} pada Eclipse~\cite{bolour2003eclipse}.

\section{\textit{Core} yang Kecil, \textit{Plugin} yang Banyak}

Susunannya terdiri atas dua bagian saja. \textbf{\textit{Core}} memuat kemampuan yang
berlaku bagi semua pemakai, misalnya menemukan bagian tambahan, menyimpan
daftarnya, dan menjalankan salah satunya atas permintaan. \textbf{\textit{Plugin}}
memuat kemampuan yang hanya dibutuhkan sebagian pemakai, misalnya dukungan satu
bahasa pemrograman tertentu.

Gambar~\ref{fig:inti-plugin} menunjukkan susunannya. Seluruh panah berangkat
dari \textit{plugin} menuju \textit{core}, dan tidak ada satu pun panah yang berangkat
dari \textit{core} menuju \textit{plugin}.

\begin{figure}[htbp]
  \centering
  \begin{tikzpicture}[
    kotak/.style={draw, rounded corners=2pt, minimum width=2.2cm,
      minimum height=0.7cm, align=center, font=\scriptsize\bfseries},
    panah/.style={-{Stealth[length=2mm]}, thick}
  ]
    \node[kotak, fill=praditagreen!25, minimum width=3.2cm,
      minimum height=1.1cm] (inti) {\\textit{Core}\\\\
      {\tiny\mdseries menemukan, mendaftar, menjalankan}};

    \node[kotak, fill=blue!8, above left=1.1cm and 0.2cm of inti] (p1)
      {\textit{plugin} konversi};
    \node[kotak, fill=blue!8, above right=1.1cm and 0.2cm of inti] (p2)
      {\textit{plugin} kurs};
    \node[kotak, fill=blue!8, below left=1.1cm and 0.2cm of inti] (p3)
      {\textit{plugin} riwayat};
    \node[kotak, fill=green!10, below right=1.1cm and 0.2cm of inti] (p4)
      {\textit{plugin} baru};

    \draw[panah] (p1) -- (inti);
    \draw[panah] (p2) -- (inti);
    \draw[panah] (p3) -- (inti);
    \draw[panah] (p4) -- (inti);
  \end{tikzpicture}
  \caption{\textit{Core} beserta keempat \textit{plugin}. Setiap panah
    berangkat dari \textit{plugin} menuju \textit{core}, tidak pernah
    sebaliknya, sehingga menambah kotak baru tidak menyentuh isi
    \textit{core}}
  \label{fig:inti-plugin}
\end{figure}

Analogi yang mendekati adalah kantin kampus yang menyewakan kios. Pengelola
menyediakan tempat, listrik, air, dan aturan jam buka, lalu siapa saja boleh
membuka kios selama aturan itu dipenuhi. Pengelola tidak perlu tahu menu apa
yang akan dijual, sebab yang diurusnya hanya daftar penyewa beserta papan
namanya. Satu kios yang tutup pun tidak menutup kantinnya, sebab pembeli cukup
berpindah ke kios sebelah.

Pembagian tersebut menentukan letak setiap keputusan.
Tabel~\ref{tab:inti-plugin} merangkum tanggung jawab dan larangan kedua bagian.

\begin{table}[htbp]
  \centering
  \caption{Tanggung jawab dan larangan \textit{core} serta \textit{plugin}}
  \label{tab:inti-plugin}
  \begin{tabularx}{\textwidth}{|l|X|X|}
    \hline
    \textbf{Bagian} & \textbf{Tanggung jawab} & \textbf{Tidak boleh mengetahui} \\
    \hline
    \textit{Core} & Menemukan \textit{manifest}, menyimpan \textit{registry}, menjalankan perintah & Nama, jumlah, dan isi \textit{plugin} mana pun \\
    \hline
    \textit{Plugin} & Menyediakan satu kemampuan beserta \textit{manifest}-nya & Isi \textit{plugin} lain, dan cara \textit{core} menyimpan \textit{registry}-nya \\
    \hline
    \textit{Contract} & Menyatakan bentuk yang wajib dipenuhi kedua belah pihak & Teknologi penyimpanan maupun bentuk antarmuka \\
    \hline
  \end{tabularx}
\end{table}

\section{\textit{Contract}, \textit{Manifest}, dan \textit{Registry}}

Tiga hal dibutuhkan agar susunan tersebut bekerja. \textbf{\textit{Contract}} menyatakan
bentuk yang wajib dipenuhi setiap \textit{plugin}. \textbf{\textit{Manifest}} menyatakan
kemampuan apa yang disumbangkan, dan disimpan di luar kode. \textbf{\textit{Registry}}
menyimpan \textit{manifest} yang sudah terbaca, sehingga \textit{core} dapat mencari tanpa
memindai ulang.

Eclipse menyebut \textit{contract} tersebut \textit{extension point}, yaitu slot yang
disediakan sebuah \textit{plugin} agar \textit{plugin} lain dapat
memenuhinya~\cite{bolour2003eclipse}. Isi \textit{manifest} dibaca lebih dahulu, lalu
tersedia lewat sebuah \textit{registry}, dan cara itulah yang membuat \textit{core} dapat
menemukan penyumbang tanpa mengimpor kelasnya.

Listing~\ref{lst:kontrak} memperlihatkan \textit{contract} pada contoh bab ini. Isinya
hanya satu bentuk pemanggilan, sebab \textit{contract} yang gemuk menyulitkan penulis
\textit{plugin}.

\begin{lstlisting}[style=PythonStyle, caption={Kontrak yang wajib dipenuhi
  setiap \textit{plugin}, tersedia di \berkas{sources/chapter05/kontrak.py}},
  label=lst:kontrak]
# Kunci yang wajib ada pada setiap file manifest.
KUNCI_MANIFES = ("nama", "perintah", "keterangan", "modul", "kelas")


class Plugin(Protocol):
  """Satu kemampuan yang disumbangkan sebuah plugin kepada core."""

  def __init__(self, inti: object) -> None:
    """Menerima core sebagai satu-satunya ketergantungan sebuah plugin."""

  def jalankan(self, argumen: list[str]) -> str:
    """Menjalankan kemampuannya lalu mengembalikan teks siap tampil."""
\end{lstlisting}

\textit{Manifest} sengaja disimpan sebagai \textit{file} JavaScript Object
Notation (JSON) yang terpisah dari modulnya. Pemisahan itu bukan soal selera.
\textit{Manifest} yang terpisah dapat dibaca tanpa menjalankan kode pluginnya,
dan kemampuan itulah yang dipakai pada seksi berikutnya. Visual Studio Code
menyimpannya di \texttt{package.json}, Chrome dan Firefox di
\texttt{manifest.json}~\cite{chrome2026content,mozilla2026manifest}, sedangkan
Eclipse dahulu memakai \texttt{plugin.xml}~\cite{bolour2003eclipse}.

Listing~\ref{lst:manifest} memperlihatkan satu \textit{manifest} utuh pada
contoh bab ini. Kelima kuncinya persis yang disebut \texttt{KUNCI\_MANIFES},
dan dua kunci terakhirlah yang menghubungkan \textit{manifest} dengan kodenya.

\begin{lstlisting}[language=json, caption={\textit{Manifest} \textit{plugin}
  konversi, tersedia di \berkas{sources/chapter05/plugins/plugin_konversi.json}},
  label=lst:manifest]
{
  "nama": "Konversi Mata Uang",
  "perintah": "konversi",
  "keterangan": "Menghitung konversi satu nominal antar mata uang",
  "modul": "plugins.plugin_konversi",
  "kelas": "ConversionPlugin"
}
\end{lstlisting}

Ketiga kunci pertama dapat dibaca tanpa menyentuh kode pluginnya sama sekali,
sedangkan dua kunci terakhir barulah yang menunjuk kodenya.
Tabel~\ref{tab:kunci-manifest} merinci peran masing-masing beserta saat kunci
itu dipakai.

\begin{table}[htbp]
  \centering
  \caption{Peran setiap kunci \textit{manifest} pada saat \textit{plugin} dimuat}
  \label{tab:kunci-manifest}
  \begin{tabularx}{\textwidth}{|l|l|X|}
    \hline
    \textbf{Kunci} & \textbf{Dipakai saat} & \textbf{Perannya} \\
    \hline
    \texttt{nama} & Pemuatan gagal & Muncul pada pesan \texttt{PluginError}, misalnya \texttt{Plugin Rusak gagal dimuat} \\
    \hline
    \texttt{perintah} & Aplikasi mulai & Kunci pada \textit{registry}, lalu dicocokkan dengan argumen baris perintah \\
    \hline
    \texttt{keterangan} & Aplikasi mulai & Dicetak pada daftar perintah, tanpa mengimpor modul \textit{plugin} \\
    \hline
    \texttt{modul} & Perintah dipanggil & Diserahkan ke \texttt{importlib.import\_module}, dan impornya terjadi di sini \\
    \hline
    \texttt{kelas} & Perintah dipanggil & Diambil lewat \texttt{getattr}, lalu objeknya dibuat dengan \texttt{Core} sebagai argumen \\
    \hline
  \end{tabularx}
\end{table}

Gambar~\ref{fig:kelas-microkernel} menunjukkan hubungan ketiganya sebagai
\textit{class diagram}. Perhatikan bahwa panah dari \textit{core} hanya menuju \textit{contract} dan
\textit{manifest}, tidak pernah menuju modul \textit{plugin}.

\begin{figure}[htbp]
  \centering
  \includegraphics[width=0.78\textwidth, height=0.52\textheight,
    keepaspectratio]{kelas-microkernel.pdf}
  \caption{\textit{Core} mengenal \textit{contract} dan \textit{manifest} saja. Ketiga \textit{plugin}
    memenuhi \textit{contract} yang sama, dan manifeslah yang menunjuk nama modulnya}
  \label{fig:kelas-microkernel}
\end{figure}

\section{Kapan \textit{Plugin} Diaktifkan}

Menemukan \textit{plugin} dan mengaktifkannya adalah dua langkah yang berbeda,
dan jarak antara keduanya menentukan \textit{startup time} sebuah aplikasi.
\textbf{\textit{Eager activation}} mengimpor seluruh modul saat aplikasi mulai.
\textbf{\textit{Lazy activation}} menunda impor sampai kemampuannya benar-benar dipakai.

Visual Studio Code memilih cara kedua. Setiap ekstensi menyatakan
\textit{activation event} pada \textit{manifest}-nya, misalnya \texttt{onLanguage} atau
\texttt{onCommand}, sehingga ekstensi baru dimuat ketika peristiwa itu terjadi.
Dokumentasinya bahkan menandai peristiwa yang menyala di segala keadaan sebagai
pilihan terakhir, sebab pilihan itu memperlambat awal
mula~\cite{vscode2026activation}. Eclipse menempuh jalan yang sama lewat
pemuatan yang ditunda, dan objek \textit{plugin} baru dibuat saat kemampuannya
dipakai~\cite{bolour2003eclipse}.

Mekanisme pemuatannya sudah lama dipakai di luar Python. Contoh Java pada
\berkas{sources/chapter05/java-example/} memuat setiap \textit{plugin} dari
\textit{file} JAR tersendiri memakai \texttt{URLClassLoader}, membaca kunci
\texttt{mainclass} pada \texttt{plugin.properties}, lalu membentuk objeknya
lewat \texttt{Class.forName} beserta konstruktor hasil refleksi.
Tabel~\ref{tab:java-python} memetakan kedua sisinya.

\begin{table}[htbp]
  \centering
  \caption{Padanan mekanisme pemuatan pada contoh Java dan contoh Python}
  \label{tab:java-python}
  \begin{tabularx}{\textwidth}{|X|X|}
    \hline
    \textbf{Contoh Java} & \textbf{Contoh Python} \\
    \hline
    \texttt{plugin.properties} berisi kunci \texttt{mainclass} & \texttt{plugins/plugin\_*.json} berisi kunci \texttt{modul} dan \texttt{kelas} \\
    \hline
    \texttt{URLClassLoader} beserta \texttt{Class.forName} & \texttt{importlib.import\_module} beserta \texttt{getattr} \\
    \hline
    Konstruktor hasil refleksi menerima \texttt{Greeter} & Konstruktor kelas menerima objek \texttt{Core} \\
    \hline
    Antarmuka \texttt{Plugin} & \textit{Protocol} \texttt{Plugin} pada \berkas{sources/chapter05/kontrak.py} \\
    \hline
  \end{tabularx}
\end{table}

Gambar~\ref{fig:sekuens-aktivasi} menunjukkan urutannya pada contoh bab ini.
\textit{Manifest} dibaca saat mulai, sedangkan impor modulnya menunggu pemanggilan
pertama.

\begin{figure}[htbp]
  \centering
  \includegraphics[width=0.95\textwidth, height=0.55\textheight,
    keepaspectratio]{sekuens-aktivasi.pdf}
  \caption{\textit{Lazy activation}. Impor hanya terjadi pada pemanggilan pertama, dan
    \textit{plugin} yang tidak pernah dipanggil tidak pernah diimpor}
  \label{fig:sekuens-aktivasi}
\end{figure}

Listing~\ref{lst:aktifkan} memperlihatkan \textit{core}-nya. Hasil impor disimpan,
sehingga pemanggilan kedua tidak membayar biaya yang sama.

\begin{lstlisting}[style=PythonStyle, caption={Aktivasi satu \textit{plugin},
  tersedia di \berkas{sources/chapter05/inti.py}}, label=lst:aktifkan]
def aktifkan(self, perintah: str) -> kontrak.Plugin:
  """Mengimpor modul, mengambil kelasnya, lalu membuat objek pluginnya."""
  if perintah in self.plugin_aktif:
    return self.plugin_aktif[perintah]  # objek pemenuh kontrak.Plugin
  manifes = self.manifes[perintah]
  try:
    modul = importlib.import_module(manifes["modul"])
    kelas = getattr(modul, manifes["kelas"])
    plugin = kelas(self)
  except Exception as kesalahan:
    self.galat_muat[perintah] = str(kesalahan)
    raise PluginError(f"{manifes['nama']} gagal dimuat: {kesalahan}")
  self.plugin_aktif[perintah] = plugin
  return plugin
\end{lstlisting}

Pembandingnya ada pada Listing~\ref{lst:aktifkan-semua}. Metode tersebut
memanggil \texttt{aktifkan} untuk setiap \textit{manifest} yang terbaca, tanpa
menunggu satu perintah pun dipanggil, sehingga keempat modul diimpor saat
aplikasi mulai. Latihan~3 mengukur harganya.

\begin{lstlisting}[style=PythonStyle, caption={Aktivasi seluruh \textit{plugin}
  sekaligus, tersedia di \berkas{sources/chapter05/inti.py}},
  label=lst:aktifkan-semua]
def aktifkan_semua(self) -> int:
  """Mengaktifkan seluruh plugin sekaligus, seperti eager activation."""
  jumlah = 0
  for perintah in self.manifes:
    try:
      self.aktifkan(perintah)
      jumlah += 1
    except PluginError:
      pass
  return jumlah
\end{lstlisting}

\section{Ketika Satu \textit{Plugin} Rusak}

Kemampuan menerima kode dari luar membawa persoalan yang tidak dimiliki
arsitektur pada bab sebelumnya. Kode itu ditulis pihak lain, kualitasnya tidak
dapat dipastikan, dan sebagiannya pasti akan gagal. Pertanyaannya, apakah
kegagalan satu \textit{plugin} ikut menjatuhkan \textit{core}-nya.

Jawabannya ditentukan oleh dua keputusan. Keputusan pertama adalah kapan modul
diimpor, sebab kegagalan saat impor jauh lebih berbahaya daripada kegagalan saat
pemanggilan. Keputusan kedua adalah apakah pemanggilannya dilindungi
\textit{error handler}. Chrome menempuh pemisahan yang lebih jauh lagi, yaitu menjalankan
\textit{content script} di dalam \textit{isolated world} yang terpisah dari
halaman maupun dari ekstensi lain~\cite{chrome2026content}.

Kasus berikut memperlihatkan akibatnya. Sebuah tim memasang lima
\textit{plugin} pada aplikasi internalnya, dan kelimanya diimpor begitu aplikasi
mulai. Satu \textit{plugin} membaca \textit{file} pengaturan di baris teratas
modulnya, sehingga \textit{file} itu dibaca sewaktu modulnya diimpor. Suatu pagi
\textit{file} pengaturan itu terhapus, sehingga impornya gagal. Impor tersebut
berjalan sebelum layar muncul, dan tidak ada \textit{error handler} di depannya,
sehingga aplikasi berhenti sebelum menampilkan apa pun. Empat \textit{plugin}
lain yang sehat ikut tidak berjalan. Keempatnya sebenarnya tidak rusak, tetapi
kelima \textit{plugin} berbagi satu proses dan satu urutan mulai, sehingga
\textit{error} yang tidak tertangkap menghentikan proses itu sebelum keempatnya
sempat berjalan. Yang rusak pun tidak dapat dinonaktifkan, sebab menu
pengaturannya tidak pernah terbuka.

Gambar~\ref{fig:proses-berhenti} memperlihatkan sebabnya.

\begin{figure}[htbp]
  \centering
  \begin{tikzpicture}[
    modul/.style={draw, rounded corners=2pt, fill=blue!8, minimum height=9mm,
      minimum width=17mm, font=\scriptsize, align=center},
    terimpor/.style={modul, fill=green!10},
    gagal/.style={modul, fill=red!10, draw=red!70, dashed},
    berhenti/.style={draw=red!70, dashed, rounded corners=2pt, fill=red!10,
      minimum height=8mm, font=\scriptsize, align=center},
    panah/.style={-{Stealth[length=2mm]}, thick},
  ]
    \node[terimpor] (p1) {konversi\\{\tiny\mdseries terimpor}};
    \node[terimpor, right=4mm of p1] (p2) {kurs\\{\tiny\mdseries terimpor}};
    \node[terimpor, right=4mm of p2] (p3) {riwayat\\{\tiny\mdseries terimpor}};
    \node[gagal, right=4mm of p3] (p4) {rusak\\{\tiny\mdseries KeyError}};
    \node[modul, right=4mm of p4] (p5) {pajak\\{\tiny\mdseries tidak terimpor}};
  
    \draw[panah] (p1) -- (p2);
    \draw[panah] (p2) -- (p3);
    \draw[panah] (p3) -- (p4);
    \draw[panah, red!70, dashed] (p4) -- (p5);
  
    \node[draw=praditagreen, rounded corners=3pt, inner sep=3.5mm,
      fit=(p1)(p5), label={[font=\scriptsize]above:Satu proses Python, satu
      urutan impor}] (proses) {};
  
    \node[berhenti, below=9mm of p4] (stop) {Proses berhenti sebelum layar
      muncul};
    \draw[panah, red!70, dashed] (p4) -- (stop)
      node[midway, right, font=\tiny] {error tidak tertangkap};
  \end{tikzpicture}
  \caption{Satu \textit{error} impor yang tidak tertangkap menghentikan seluruh proses. \textit{Plugin} yang sudah terimpor ikut hilang, sedangkan yang belum tidak pernah sempat diimpor}
  \label{fig:proses-berhenti}
\end{figure}

Tabel~\ref{tab:isolasi} merangkum kedua jalur tersebut beserta akibatnya.

\begin{table}[htbp]
  \centering
  \caption{Nasib aplikasi ketika satu \textit{plugin} gagal diimpor}
  \label{tab:isolasi}
  \begin{tabularx}{\textwidth}{|l|X|X|}
    \hline
    \textbf{Jalur} & \textbf{Yang terjadi} & \textbf{Akibat bagi pengguna} \\
    \hline
    Tanpa \textit{error handler} & \textit{Error} naik sampai ke pemanggil paling luar & Aplikasi berhenti, \textit{plugin} sehat ikut berhenti \\
    \hline
    Dengan \textit{error handler} & \textit{Error} dicatat, perintah lain tetap dilayani & Satu kemampuan hilang, sisanya tetap jalan \\
    \hline
    Proses terpisah & \textit{Error} tertahan di dalam prosesnya sendiri & Aplikasi utuh, tetapi biayanya lebih mahal \\
    \hline
  \end{tabularx}
\end{table}

\section{\textit{Plugin} pada Produk Nyata}

Keempat produk berikut memakai susunan yang sama dengan nama yang berbeda.
Tabel~\ref{tab:produk-plugin} merangkum istilah dan \textit{file} \textit{manifest} masing-masing.

\begin{table}[htbp]
  \centering
  \caption{Sebutan dan \textit{file} \textit{manifest} pada empat produk yang memakai
    \textit{plugin}}
  \label{tab:produk-plugin}
  \begin{tabularx}{\textwidth}{|l|l|X|}
    \hline
    \textbf{Produk} & \textbf{\textit{Manifest}} & \textbf{Catatan} \\
    \hline
    Eclipse & \texttt{plugin.xml} & Kontraknya disebut \textit{extension point}, dan pemenuhnya disebut \textit{extension} \\
    \hline
    Visual Studio Code & \texttt{package.json} & Pemuatan ditunda sampai \textit{activation event} yang didaftarkan terjadi \\
    \hline
    Chrome & \texttt{manifest.json} & \textit{Content script} berjalan di \textit{isolated world} yang terpisah \\
    \hline
    Firefox & \texttt{manifest.json} & Memakai model WebExtensions yang sama, dengan sejumlah perbedaan yang disengaja \\
    \hline
  \end{tabularx}
\end{table}

Perbedaan namanya menutupi kesamaan susunannya. Ketiganya menyimpan \textit{manifest} di
luar kode, menunda pemuatan, dan menolak menyebut nama penyumbang di dalam
\textit{core}-nya. Yang berbeda hanyalah seberapa jauh pemisahannya dibawa, dan Chrome
membawanya paling jauh dengan memisahkan lingkungan eksekusinya.

\section{Ringkasan}

Microkernel membagi sistem menjadi \textit{core} yang kecil beserta bagian tambahan yang
dapat dipasang dan dicabut. Bentuk yang paling sering ditemui bukan kernel
sistem operasi, melainkan arsitektur \textit{plugin} pada perkakas sehari-hari.
Eclipse, Visual Studio Code, Chrome, dan Firefox sama-sama memakainya.

Yang membuat susunan tersebut bekerja adalah arah pengenalan. \textit{Core} menulis
\textit{contract} dan membaca \textit{manifest}, sedangkan \textit{plugin} yang
memenuhi \textit{contract} tersebut ditulis pihak lain. Karena \textit{core}
tidak pernah menyebut nama satu penyumbang pun, menambah kemampuan berarti
menambah \textit{file}, bukan menyunting \textit{core}. Latihan~1 membuktikan
klaim itu dengan menghitung penyebutannya, dan angkanya nol.

Karena \textit{manifest} berupa \textit{file} terpisah dari kodenya, daftar perintah
dapat dibaca tanpa mengimpor satu modul \textit{plugin} pun. Pemisahan itu
menimbulkan pilihan kedua, yaitu kapan modulnya diimpor. \textit{Lazy
activation} membaca \textit{manifest} saat mulai lalu menunda impor sampai
perintahnya dipanggil. Pada contoh bab ini, cara tersebut
membuat \textit{startup time} 3,4 kali lebih cepat, dan selisihnya lebih besar daripada
\textit{range} pengukurannya sendiri. Kesimpulan itu berbeda dengan Bab~2, dan
perbedaannya justru memperlihatkan gunanya membandingkan selisih dengan
\textit{range}.

Kemampuan menerima kode dari luar menuntut harga tersendiri. Satu
\textit{plugin} yang gagal saat impor dapat menjatuhkan seluruh aplikasi bila
pemanggilannya tidak dilindungi \textit{error handler}. Latihan~2 memperlihatkan
selisihnya, yaitu tiga perintah tetap berjalan pada jalur dengan \textit{error
handler}, sedangkan jalur tanpa \textit{error handler} berhenti pada
\textit{file} yang rusak.

Pesan utama bab ini karena itu bukan bahwa \textit{core} harus kecil, melainkan bahwa
\textit{core} harus buta terhadap penyumbangnya. Kebutaan itu yang dibayar dengan
\textit{contract}, \textit{manifest}, dan \textit{registry}, dan ketiganya baru terasa mahal ketika jumlah
\textit{plugin} masih sedikit.

\section{Latihan}

Ketiga latihan berikut memakai satu basis kode yang sama di
\berkas{sources/chapter05/}, tetapi menyasar tiga masalah yang berbeda.
Latihan~1 memeriksa \textit{core independence}. Latihan~2
memeriksa akibat satu \textit{plugin} yang rusak. Latihan~3 mengukur selisih
\textit{startup time} antara kedua cara aktivasi. Hasil setiap mahasiswa dapat berbeda,
dan yang dilaporkan adalah pola perbandingannya, bukan angka mutlaknya.

Venv yang sama dipakai kembali, dan pengaktifannya seperti pada
Listing~\ref{lst:venv-ch05}. Latihan bab ini hanya memakai pustaka standar
Python, sehingga tidak ada pemasangan tambahan.

\begin{lstlisting}[language=bash, caption={Mengaktifkan venv sebelum ketiga
  latihan, skrip tersedia di \berkas{sources/siapkan_venv.sh}},
  label=lst:venv-ch05]
bash sources/siapkan_venv.sh
source .venv/bin/activate
cd sources/chapter05
python inti.py daftar
# Keluarannya baris pertama: Plugin terpasang: 4
\end{lstlisting}

\subsection*{Latihan 1: \textit{Core Independence}}

Latihan ini menguji klaim utama arsitektur \textit{plugin}, yaitu menambah
kemampuan cukup dengan menambah \textit{file}. Klaim tersebut menyebut angka, sehingga
dapat dibuktikan salah. Satu penyebutan nama modul \textit{plugin} di dalam \textit{core}
sudah cukup untuk menggugurkannya. Langkah kerjanya sebagai berikut.

\begin{enumerate}
  \item Jalankan pemeriksaan seperti pada Listing~\ref{lst:jalankan-periksa-inti},
    lalu catat baris \texttt{Plugin disebut core} beserta ketiga angka di
    bawahnya.
  \item Buka \berkas{sources/chapter05/inti.py}, lalu cari nama modul
    \textit{plugin} di dalamnya. Catat hasil pencarian itu.
  \item Salin \berkas{sources/chapter05/plugins/plugin_riwayat.py} beserta
    \textit{manifest}-nya menjadi \texttt{plugin\_pajak.py}, lalu ubah nama, perintah, dan
    nama kelasnya.
  \item Jalankan \texttt{python inti.py daftar}, lalu catat jumlah
    \textit{plugin} yang terbaca.
  \item Jalankan ulang \texttt{python periksa\_inti.py}, lalu catat berapa baris
    \textit{core} yang berubah selama langkah ketiga dan keempat.
\end{enumerate}

\begin{lstlisting}[language=bash, caption={Menghitung penyebutan \textit{plugin}
  di dalam inti}, label=lst:jalankan-periksa-inti]
python periksa_inti.py
# Keluarannya: Plugin disebut core    : 0 []
\end{lstlisting}

Tabel~\ref{tab:lat1plugin-contoh} berisi hasil satu kali pemeriksaan nyata.

\begin{table}[htbp]
  \centering
  \caption{Contoh hasil Latihan 1 yang sudah terisi}
  \label{tab:lat1plugin-contoh}
  \begin{tabularx}{\textwidth}{|X|l|}
    \hline
    \textbf{Yang dilaporkan} & \textbf{Nilai} \\
    \hline
    \textit{Manifest} ditemukan & 4 \\
    \hline
    \textit{Plugin} disebut \textit{core} & 0 \\
    \hline
    Baris \textit{core} & 142 \\
    \hline
    Baris seluruh \textit{plugin} & 112 \\
    \hline
    Bagian \textit{plugin} & 44 persen \\
    \hline
  \end{tabularx}
\end{table}

Isikan hasil pemeriksaan pada Tabel~\ref{tab:lat1plugin-hasil}.

\begin{table}[htbp]
  \centering
  \caption{Lembar isian Latihan 1}
  \label{tab:lat1plugin-hasil}
  \begin{tabularx}{\textwidth}{|X|l|l|}
    \hline
    \textbf{Yang dilaporkan} & \textbf{Sebelum} & \textbf{Sesudah} \\
    \hline
    \textit{Manifest} ditemukan & \dots & \dots \\
    \hline
    \textit{Plugin} disebut \textit{core} & \dots & \dots \\
    \hline
    Baris \textit{core} & \dots & \dots \\
    \hline
    Baris seluruh \textit{plugin} & \dots & \dots \\
    \hline
  \end{tabularx}
\end{table}

Jawab tiga pertanyaan berikut dengan menunjuk angka pada tabel. Pertama, baris
\textit{core} tidak berubah sama sekali walaupun satu \textit{plugin} ditambahkan.
Sebutkan bagian \textit{core} yang membuat hal itu mungkin. Kedua, jelaskan apa yang akan
terjadi pada angka penyebutan bila \textit{core} mengimpor satu \textit{plugin} secara
langsung demi kemudahan. Ketiga, hitung bagian kode yang berupa \textit{plugin},
lalu bandingkan dengan perkiraan awal sebelum pengukuran dilakukan.

\subsection*{Latihan 2: \textit{Error Isolation}}

Latihan ini membuktikan bahwa kegagalan satu \textit{plugin} tidak harus
menjatuhkan \textit{core}-nya. \textit{File} \berkas{sources/chapter05/plugins/plugin_rusak.py} sengaja
gagal pada tingkat modul, sehingga kegagalannya muncul saat impor. Langkah
kerjanya sebagai berikut.

\begin{enumerate}
  \item Baca \berkas{sources/chapter05/plugins/plugin_rusak.py}, lalu
    tunjukkan baris yang menyebabkan kegagalannya.
  \item Jalankan ketiga perintah pada Listing~\ref{lst:jalankan-uji-isolasi},
    lalu catat kedua baris hasil \texttt{uji\_isolasi.py}.
  \item Perhatikan dua perintah terakhir pada listing tersebut, lalu catat
    apakah perintah kedua tetap dilayani sesudah perintah pertama gagal.
  \item Buka \berkas{sources/chapter05/inti.py}, lalu temukan \textit{error
    handler} yang menangkap kegagalan impor.
  \item Hapus sementara \textit{error handler} tersebut, jalankan ulang
    \texttt{python inti.py rusak}, lalu catat akibatnya. Perintah
    \texttt{kurs} sengaja tidak dipakai di sini, sebab \textit{plugin} sehat
    tidak pernah menyentuh \textit{error handler} itu. Kembalikan kodenya
    sesudah dicatat.
\end{enumerate}

\begin{lstlisting}[language=bash, caption={Menjalankan satu \textit{plugin} yang
  rusak beserta perintah sesudahnya}, label=lst:jalankan-uji-isolasi]
python uji_isolasi.py
# Keluarannya: Dengan error handler : 3 perintah jalan, 1 dilewati
#              Tanpa error handler  : 3 modul terimpor, lalu berhenti
python inti.py rusak
# Keluarannya: Dilewati: Plugin Rusak gagal dimuat: ('JPY', 'IDR')
python inti.py kurs
# Keluarannya: EUR -> IDR: 17,600.00
#              SGD -> IDR: 12,100.00
#              USD -> IDR: 16,250.00
\end{lstlisting}

Tabel~\ref{tab:lat2plugin-contoh} berisi hasil satu kali percobaan nyata.

\begin{table}[htbp]
  \centering
  \caption{Contoh hasil Latihan 2 yang sudah terisi}
  \label{tab:lat2plugin-contoh}
  \begin{tabularx}{\textwidth}{|l|l|X|}
    \hline
    \textbf{Jalur} & \textbf{Perintah jalan} & \textbf{Catatan} \\
    \hline
    Dengan \textit{error handler} & 3 dari 4 & Satu \textit{plugin} dilewati, sisanya dilayani \\
    \hline
    Tanpa \textit{error handler} & 3 modul & Berhenti pada \texttt{plugin\_rusak} \\
    \hline
    \textit{Error handler} dihapus & 0 dari 1 & \texttt{python inti.py rusak} berhenti dengan \texttt{KeyError} \\
    \hline
  \end{tabularx}
\end{table}

Isikan hasil percobaan pada Tabel~\ref{tab:lat2plugin-hasil}.

\begin{table}[htbp]
  \centering
  \caption{Lembar isian Latihan 2}
  \label{tab:lat2plugin-hasil}
  \begin{tabularx}{\textwidth}{|l|l|X|}
    \hline
    \textbf{Jalur} & \textbf{Perintah jalan} & \textbf{Catatan} \\
    \hline
    Dengan \textit{error handler} & \dots & \dots \\
    \hline
    Tanpa \textit{error handler} & \dots & \dots \\
    \hline
    \textit{Error handler} dihapus & \dots & \dots \\
    \hline
  \end{tabularx}
\end{table}

Jawab tiga pertanyaan berikut. Pertama, jelaskan mengapa kegagalan saat impor
lebih berbahaya daripada kegagalan saat pemanggilan. Kedua, sebutkan apa yang
hilang bagi pengguna ketika satu \textit{plugin} dilewati, lalu bandingkan
dengan yang hilang ketika aplikasi tidak mau mulai. Ketiga, jelaskan mengapa
pemisahan proses seperti pada Chrome lebih kuat daripada \textit{error handler},
beserta harga yang dibayarnya.

\subsection*{Latihan 3: \textit{Activation Cost}}

Latihan ini mengukur selisih \textit{startup time} antara \textit{lazy activation} dan \textit{eager activation}. Pengukurannya diulang seratus putaran, sebab satu putaran tidak cukup
untuk menyimpulkan apa pun. \textit{Import cache} Python dibersihkan pada setiap putaran,
agar kedua jalur menghadapi keadaan yang sama. Langkah kerjanya sebagai berikut.

\begin{enumerate}
  \item Jalankan pengukuran seperti pada Listing~\ref{lst:jalankan-ukur-aktivasi}.
    Skrip tersebut mencetak baris kemajuan setiap sepuluh putaran, sehingga
    jalannya terlihat.
  \item Catat kedelapan angka pada laporan akhirnya.
  \item Buka \berkas{sources/chapter05/ukur_aktivasi.py}, lalu tunjukkan baris
    yang membersihkan \textit{import cache}.
  \item Hapus sementara pembersihan tersebut, jalankan ulang, lalu catat
    perubahan angkanya. Kembalikan kodenya sesudah dicatat.
\end{enumerate}

\begin{lstlisting}[language=bash, caption={Mengukur waktu siap kedua cara
  aktivasi, seratus putaran}, label=lst:jalankan-ukur-aktivasi]
python ukur_aktivasi.py
# Keluarannya: Putaran lazy lebih cepat     : 100 dari 100
\end{lstlisting}

Tabel~\ref{tab:lat3plugin-contoh} berisi hasil satu kali pengukuran nyata.

\begin{table}[htbp]
  \centering
  \caption{Contoh hasil Latihan 3 yang sudah terisi, seratus putaran}
  \label{tab:lat3plugin-contoh}
  \begin{tabularx}{\textwidth}{|X|l|}
    \hline
    \textbf{Yang dilaporkan} & \textbf{Nilai} \\
    \hline
    Median \textit{lazy activation} & 0,3988 ms \\
    \hline
    Median \textit{eager activation} & 1,3380 ms \\
    \hline
    Selisih median antar jalur & 0,9391 ms \\
    \hline
    \textit{Eager activation} lebih lambat & 3,4 kali \\
    \hline
    \textit{Range} di dalam jalur \textit{lazy} & 0,2579 ms \\
    \hline
    \textit{Range} dibagi selisih & 0,3 kali \\
    \hline
    Putaran \textit{lazy} lebih cepat & 100 dari 100 \\
    \hline
  \end{tabularx}
\end{table}

Baris keenam itulah yang membedakan bab ini dari Bab~2. Di sana \textit{range}
di dalam satu jalur puluhan kali lebih besar daripada selisih yang dicari,
sehingga selisihnya belum dapat disebut ada. Di sini \textit{range}-nya justru
sepertiga selisihnya, dan pencacahan per putaran menegaskannya dengan angka
100 dari 100. Selisih yang seperti inilah yang layak dilaporkan.

Isikan hasil pengukuran pada Tabel~\ref{tab:lat3plugin-hasil}.

\begin{table}[htbp]
  \centering
  \caption{Lembar isian Latihan 3}
  \label{tab:lat3plugin-hasil}
  \begin{tabularx}{\textwidth}{|X|l|l|}
    \hline
    \textbf{Yang dilaporkan} & \textbf{Dengan pembersihan} & \textbf{Tanpa pembersihan} \\
    \hline
    Median \textit{lazy activation} & \dots & \dots \\
    \hline
    Median \textit{eager activation} & \dots & \dots \\
    \hline
    Selisih median antar jalur & \dots & \dots \\
    \hline
    \textit{Range} dibagi selisih & \dots & \dots \\
    \hline
    Putaran \textit{lazy} lebih cepat & \dots & \dots \\
    \hline
  \end{tabularx}
\end{table}

Jawab tiga pertanyaan berikut dengan menunjuk angka pada tabel. Pertama,
jelaskan mengapa selisih pada bab ini layak disebut ada, sedangkan selisih pada
Bab~2 tidak. Kedua, jelaskan apa yang terjadi pada angka ketika pembersihan
cache dihapus, lalu sebutkan pengukuran mana yang menjadi tidak jujur karenanya.
Ketiga, dengan menggabungkan hasil ketiga latihan, tentukan kapan \textit{eager activation}
tetap layak dipilih walaupun lebih lambat, lalu sebutkan angka yang menjadi
dasar keputusannya.
